\section{Conclusion and Future Work}\label{sec:concl}

The Boolean-kernel basis $\{K_b\}$ gives an exact dictionary between two operations that are usually treated separately.
One is even--odd folding of univariate polynomials over smooth multiplicative domains; the other is restriction of a variable of a multilinear polynomial in table form.
In kernel coordinates the normalised fold $\mathrm{pfold}_r$ \emph{is} the restriction (\cref{thm:fold-dictionary}).
A Reed--Solomon codeword can therefore carry a table directly, and the BaseFold paradigm runs on it with no change of representation.
We gave a self-contained proof of soundness, evaluation binding and round-by-round knowledge soundness in the unique-decoding regime, whose only external input about codes is the correlated-agreement theorem for Reed--Solomon codes (with efficient unique decoding for the extractor).
With the BCS theorem for interactive oracle reductions of~\cite{CDHZ25}, this makes the non-interactive scheme knowledge sound, for a single opening, against quantum adversaries in the quantum random oracle model (\cref{cor:pq}).
We also gave an implementation that follows the compiler of~\cite{CDHZ25}, with classical and post-quantum parameters, in which the kernel encoding and coefficient-form folding are indistinguishable in running time at the resolution of our machine.

\subsection{Future work}\label{sec:concl:future}

\paragraph{Soundness.}
\begin{itemize}
  \item \emph{List decoding.} Extend \cref{thm:sound,thm:rbr} to proximity parameters up to the Johnson bound, in the spirit of~\cite{Hab24,Zei26}. The line form of \cref{lem:line} reduces each fold to a Reed--Solomon line, so correlated agreement beyond unique decoding~\cite{BCIKS23} is the natural tool; the main work is to replace the unique codewords $c_j$ of the passing-set argument by lists. This would reduce the number of queries, in the classical parameters of \cref{rem:params} and in the post-quantum parameters of \cref{rem:pq-params}, where the quadratic loss of \cref{thm:cdhz} makes the query count conservative.
  \item \emph{Batched and repeated openings.} Prove soundness of the batched opening of \cref{rem:batch}, by correlated agreement applied to the random linear combination of the committed words, and extend \cref{cor:pq} to several openings of one commitment.
\end{itemize}

\paragraph{Zero knowledge.}
The kernel encoding is compatible with the standard masking techniques for folding-based schemes: committing to $U_f + X^N R$ for a random polynomial $R$ of small degree, masking the sumcheck with a random polynomial in the manner of Libra~\cite{XZZPS19}, and masking the folded words with a random codeword. The task is to prove that the fold dictionary and the soundness argument survive the extra degree, and that the resulting view is simulable.

\paragraph{Implementation.}
\begin{itemize}
  \item \emph{Proof size.} Folding several variables per round, with the kernel fold applied to $2^k$-element fibres, and merging authentication paths across queries reduce the number of paths per proof; together with list decoding, these are the techniques behind the smaller proofs of WHIR in \cref{tab:whir}.
  \item \emph{Engineering.} SIMD field arithmetic and a parallel verifier. None of them interacts with the encoding.
\end{itemize}

\paragraph{Formal verification.}
Replace the axiom \texttt{bciks\_unique} by a formal proof of correlated agreement in the unique-decoding regime, state relaxed round-by-round knowledge soundness for a general interactive oracle reduction in Lean, and link the Rust implementation to the Lean specification, for instance through test vectors extracted from Lean.

\paragraph{Other domains.}
The dictionary uses a smooth multiplicative subgroup in odd characteristic. Is there a kernel-type basis whose natural fold is table-form restriction on circle domains~\cite{HLP24} or on additive subspaces~\cite{LCH14,DP24}?

\subsection{Artefacts}\label{sec:concl:artefacts}
The Rust library, the Lean files, the benchmark data and the numerical checks are available at \url{https://github.com/qoosmo/kbfold} under the MIT or Apache-2.0 licences (code) and CC~BY~4.0 (paper); this version of the paper corresponds to the release tagged \texttt{v\kbversion}, and the library is published on crates.io as \texttt{kbfold}~\kbversion.
\begin{itemize}
  \item \emph{Tests} (\cref{sec:impl:tests}): \texttt{cargo test --release} in \texttt{rust/}.
  \item \emph{Tables of \cref{sec:exp}}:\\ \texttt{cargo run --release --example bench -- scaling}, and likewise with \texttt{stop} and \texttt{rate}; the recorded data are in \texttt{rust/results/}.
  \item \emph{Numerical checks} (\cref{sec:impl:tests}):\\ \texttt{cargo run --release --example paper\_checks}.
  \item \emph{Post-quantum parameters} (\cref{rem:pq-params}):\\ \texttt{cargo run --release --example pq\_params}.
  \item \emph{Lean}: \texttt{lake build} in \texttt{lean/} (Lean~4.23.0, Mathlib \texttt{v4.23.0}); \texttt{lean/KBFold/Audit.lean} prints the axioms of the main theorems.
\end{itemize}
